\documentclass[10pt,a4paper]{article}
% Annexe technique de la note « Taxer mieux plutôt que davantage » (17/11/2026). Les annexes B à E sont écrites par
% les programmes de note8p/calculs ; A est rédigée à la main ; F par calculs/annexe_bouclage.py.
% Préambule de la note de 8 pages et de son annexe technique (note8p/).
\usepackage[T1]{fontenc}
\usepackage[utf8]{inputenc}
\usepackage{lmodern}
\usepackage[french]{babel}
\usepackage[a4paper,left=2.1cm,right=2.1cm,top=1.9cm,bottom=2.0cm]{geometry}
\usepackage{graphicx}
\usepackage{booktabs,tabularx,array}
\usepackage{xcolor}
\usepackage{titlesec}
\usepackage[labelfont=bf,labelsep=period,font=small,singlelinecheck=false,justification=raggedright,skip=3pt]{caption}
\usepackage{fancyhdr}
\usepackage{enumitem}
\usepackage{microtype}
\usepackage{textcomp}
\usepackage{amsmath}
\usepackage{wrapfig}
\providecommand{\euro}{\texteuro}
\usepackage[hidelinks]{hyperref}

\definecolor{encre}{HTML}{1A1A1A}
\definecolor{bleu}{HTML}{173F8A}
\definecolor{brique}{HTML}{B8642A}
\definecolor{gris}{HTML}{555555}

\titleformat{\section}{\normalfont\large\bfseries\color{bleu}}{}{0em}{}
\titleformat{\subsection}{\normalfont\normalsize\bfseries\color{bleu}}{}{0em}{}
\titlespacing*{\section}{0pt}{1.6ex plus .6ex minus .2ex}{0.6ex plus .2ex}
\titlespacing*{\subsection}{0pt}{1.2ex plus .5ex minus .2ex}{0.4ex plus .1ex}

\pagestyle{fancy}
\fancyhf{}
\renewcommand{\headrulewidth}{0pt}
\fancyfoot[R]{\small\color{gris}\thepage}
\fancyfoot[L]{\small\color{gris}Taxer mieux plutôt que davantage}

\setlength{\parindent}{0pt}
\setlength{\parskip}{0.45em}
\linespread{1.02}
\setlist{itemsep=1pt,topsep=1pt,parsep=0pt,leftmargin=1.2em}
\setlength{\textfloatsep}{8pt plus 2pt minus 2pt}
\setlength{\floatsep}{6pt plus 2pt minus 2pt}
\renewcommand{\topfraction}{0.85}\renewcommand{\bottomfraction}{0.7}
\renewcommand{\textfraction}{0.12}\renewcommand{\floatpagefraction}{0.8}

% \graphique{fichier}{titre-message}{label} : figure pleine largeur, source composée en petit sous l'image.
\newcommand{\graphique}[3]{%
  \begin{figure}[!htbp]\centering
  \caption{#2}\label{#3}\vspace{1pt}
  \includegraphics[width=\linewidth]{figures/#1.pdf}\par\vspace{1pt}
  {\raggedright\scriptsize\color{gris}\input{figures/#1_source.tex}\par}
  \end{figure}}
% \deuxgraphiques{f1}{titre1}{lab1}{f2}{titre2}{lab2} : deux figures demi-largeur côte à côte.
\newcommand{\deuxgraphiques}[6]{%
  \begin{figure}[!htbp]
  \begin{minipage}[t]{0.485\linewidth}\caption{#2}\label{#3}\vspace{1pt}
  \includegraphics[width=\linewidth]{figures/#1.pdf}\par
  {\raggedright\scriptsize\color{gris}\input{figures/#1_source.tex}\par}\end{minipage}\hfill
  \begin{minipage}[t]{0.485\linewidth}\caption{#5}\label{#6}\vspace{1pt}
  \includegraphics[width=\linewidth]{figures/#4.pdf}\par
  {\raggedright\scriptsize\color{gris}\input{figures/#4_source.tex}\par}\end{minipage}
  \end{figure}}

\newenvironment{encadre}[1]
  {\par\smallskip\noindent\begin{minipage}{\linewidth}\color{encre}\setlength{\parskip}{0.3em}%
   \hrule height 0.6pt\vspace{4pt}{\footnotesize\bfseries\color{bleu}\MakeUppercase{#1}}\par\small}
  {\par\vspace{3pt}\hrule height 0.3pt\end{minipage}\par\smallskip}

\DeclareUnicodeCharacter{2212}{\ensuremath{-}}
\DeclareUnicodeCharacter{2248}{\ensuremath{\approx}}
\DeclareUnicodeCharacter{00D7}{\ensuremath{\times}}
\DeclareUnicodeCharacter{2192}{\ensuremath{\rightarrow}}
\DeclareUnicodeCharacter{2264}{\ensuremath{\leq}}
\DeclareUnicodeCharacter{2265}{\ensuremath{\geq}}
\DeclareUnicodeCharacter{202F}{\,}
\DeclareUnicodeCharacter{2009}{\,}
\DeclareUnicodeCharacter{03C3}{\ensuremath{\sigma}}
\DeclareUnicodeCharacter{03C4}{\ensuremath{\tau}}
\DeclareUnicodeCharacter{03C1}{\ensuremath{\rho}}
\DeclareUnicodeCharacter{03B5}{\ensuremath{\varepsilon}}
\DeclareUnicodeCharacter{03B4}{\ensuremath{\delta}}
\DeclareUnicodeCharacter{03BB}{\ensuremath{\lambda}}
\DeclareUnicodeCharacter{00B2}{\textsuperscript{2}}

% \graphiquemarge{fichier}{titre}{label} : figure demi-largeur habillée par le texte (à placer en début de paragraphe).
\newcommand{\graphiquemarge}[3]{%
  \begin{wrapfigure}{r}{0.47\linewidth}\vspace{-10pt}
  \captionsetup{type=figure}\caption{#2}\label{#3}\vspace{1pt}
  \includegraphics[width=\linewidth]{figures/#1.pdf}\par
  {\raggedright\scriptsize\color{gris}\input{figures/#1_source.tex}\par}\vspace{-6pt}
  \end{wrapfigure}}

\fancyfoot[L]{\small\color{gris}Taxer mieux plutôt que davantage — annexe technique}
\setcounter{secnumdepth}{2}
\renewcommand{\thesection}{\Alph{section}}
\renewcommand{\thesubsection}{\thesection.\arabic{subsection}}
\titleformat{\section}{\normalfont\large\bfseries\color{bleu}}{\thesection.}{0.5em}{}
\titleformat{\subsection}{\normalfont\normalsize\bfseries\color{bleu}}{\thesubsection}{0.5em}{}

\begin{document}

{\footnotesize\color{brique}\bfseries ANNEXE TECHNIQUE\quad·\quad NOVEMBRE 2026\par}
\vspace{4pt}
{\fontsize{18}{22}\selectfont\bfseries\color{bleu}Taxer mieux plutôt que davantage\par}
\vspace{3pt}
{\large Données, hypothèses, formules, variantes et trajectoire budgétaire\par}
\vspace{5pt}
{\color{bleu}\hrule height 1pt}
\vspace{4pt}

{\small Chaque chiffre de la note est produit par un programme du dossier \texttt{fiscalite/note8p/calculs} et
enregistré dans \texttt{note8p/sorties} (fichiers \texttt{.json}, et comptes rendus \texttt{.md} qui détaillent les
sources). Statuts : \textbf{[E]} résultat empirique publié ; \textbf{[S]} simulation calibrée par nous ; \textbf{[H]}
hypothèse ; \textbf{[C]} choix de politique. Base : prélèvements obligatoires de 2025 ; droit de 2026 ; répartition de
l'impôt sur le revenu entre foyers : revenus de 2024 (IRCOM), les derniers publiés.\par}

{\setcounter{tocdepth}{1}\small\tableofcontents}

\section{Mesurer l'efficacité : la perte sèche, et ce qu'elle n'est pas}\label{ann:methode}

\textbf{L'idée.} Un salarié pourrait faire une mission de plus payée 100~€ ; il n'y tient qu'à hauteur de 80~€. Sans
impôt, il la fait : la société y gagne 20~€. Si l'on prélève 30~€, il ne lui reste que 70~€ : il renonce. L'État ne
touche rien, le salarié non plus, et les 20~€ de valeur ont disparu. C'est la perte sèche : la valeur des activités
que l'impôt fait abandonner, au-delà de ce qu'il rapporte. Une réforme est plus efficace si, à recette égale, elle fait
renoncer à moins d'activités utiles.

\textbf{La formule.} Nous utilisons la même formule pour les quatre réformes (Harberger, 1964 ; Hines, 1999) :
\[
\Delta PS \;=\; -\sum_j \big(\bar T_j - s_j\,\tau_{k(j)}\big)\,\Delta Q_j ,\qquad
\Delta Q_j = \varepsilon_j\,Q_j\,\Delta\ln(1-t_j).
\]
$j$ parcourt les marges sur lesquelles une réforme agit : heures et effort, décision de travailler, ventes de
logements, construction, stock d'équipements. $Q_j$ est la quantité concernée, $t_j$ le taux auquel elle réagit,
$\varepsilon_j$ son élasticité. $\Delta Q_j$ est la quantité supplémentaire (ou perdue) ; chaque unité vaut, pour la
collectivité, le coin total $\bar T_j$ qui la frappe, moyenne de sa valeur avant et après réforme (c'est la version
« trapèze » de la formule, exacte au second ordre).

\textbf{Les coins préexistants.} Sur une marge déjà taxée, une unité d'activité en moins coûte tout ce qui était
prélevé dessus, pas seulement l'impôt réformé. Pour le travail, le coin total comprend l'impôt fusionné, les cotisations
qui n'ouvrent pas de droits et le retrait de la prime d'activité : nous ne réformons pas ces prélèvements, mais les
ignorer sous-estimerait la valeur d'un emploi de plus au voisinage du Smic. Pour l'immobilier, le coin est le droit de
mutation sur les ventes et l'impôt sur le bâti sur la construction ; l'impôt sur le sol n'a pas de marge, puisque la
quantité de terrain ne réagit pas. Pour l'IS, c'est l'écart entre le rendement avant et après impôt de l'équipement
marginal (coût d'usage de Hall et Jorgenson).

\textbf{Le réel, pas le déclaré.} Seule la baisse réelle de l'activité entre dans le gain. Un revenu déclaré plus tard,
autrement ou dans une société est un effet sur la recette, donné à part, net de ce qui est récupéré sur une autre
assiette ($s_j$, part déplacée, taxée au taux $\tau_k$ ; Chetty, 2009). Nous ne transformons jamais une élasticité du
revenu imposable en élasticité de la production.

\textbf{Les marges et les élasticités.} Heures et effort (revenu réel) : $\varepsilon$ de 0{,}02 / 0{,}05 / 0{,}10 sous
le revenu médian et de 0{,}05 / 0{,}10 / 0{,}20 à 0{,}25 au-dessus (jeux bas, central, haut ; ordre de grandeur de
Lehmann, Marical et Rioux, 2013, sur données françaises). Décision de travailler : élasticité au taux
de prélèvement sur tout le salaire d'entrée de 0{,}10 / 0{,}25 / 0{,}50 sous deux Smic (synthèse de Chetty, Guren,
Manoli et Weber, 2013). Ventes de logements : 0 à 8~\% par point de droits, 6 au centre ; 0 si la réaction n'est que
temporaire (Bérard et Trannoy, 2018). Stock d'équipements : élasticité au coût d'usage de 0{,}5 à 1{,}6 (Maffini, Xing et
Devereux, 2019). Les pensions, l'épargne et l'âge de départ ne réagissent pas dans le calcul, faute de mesure fiable.

\textbf{La perception de la CSG.} Les salariés réagissent au taux qu'ils perçoivent. Le nouvel impôt, publié avec un
barème, est supposé perçu en entier ; la CSG d'aujourd'hui l'est avec un poids $\theta$ inconnu. Si $\theta=1$, la
fusion réduit la perte sèche de 3{,}8~Md€ ; elle ne gagne plus rien sous $\theta \approx 0{,}48$ ; à $\theta=0$, elle
la relève de 3{,}3~Md€. C'est l'hypothèse qui peut le plus changer la conclusion.

\textbf{Pas de doubles comptes.} La baisse de l'IS financée par l'excédent foncier est comptée dans le bloc IS, pas
dans le bloc foncier ; la déductibilité de la taxe foncière est comptée une fois, dans le foncier. Les effets de l'IS sur
le capital et le PIB ne sont pas ajoutés à sa perte sèche : ce sont deux mesures du même mécanisme. La perte de valeur
des terrains due à la taxe est un transfert entre propriétaires actuels et collectivité, pas une perte sèche. Aucune
perte sèche n'est chiffrée pour les plus-values : aucune élasticité ne s'applique proprement à une règle de transmission.

\textbf{Ce que la perte sèche n'est pas.} Ce n'est pas un supplément de PIB : elle mesure une valeur nette (le revenu
d'une heure de travail en plus moins la valeur du loisir sacrifié), pas une production brute. C'est une mesure statique,
de niveau, à comportements mesurés sur de petites réformes. Elle ne dit rien du taux de croissance. Nous n'en faisons
pas la somme avec les effets réels, et nous ne donnons pas de chiffre global de croissance.

\begin{table}[!htbp]\centering\small\setlength{\tabcolsep}{4pt}
\caption{Ce que l'on peut chiffrer, réforme par réforme}\label{tab:gains}
\begin{tabularx}{\linewidth}{@{}>{\raggedright\arraybackslash}p{2.6cm}>{\raggedright\arraybackslash}X>{\raggedright\arraybackslash}p{4.4cm}@{}}
\toprule
Réforme & Perte sèche évitée (Md€ par an ; central, fourchette) & Effets réels \\
\midrule
Nouvel impôt sur le revenu & 3,8 (1,9 à 9,0) si la CSG est perçue comme l'IR ; 0,1 si à moitié ; $-3{,}3$ si pas du
tout [S/H] & offre de travail : +305\,000 personnes (122\,000 à 610\,000) si perçue [S/H] \\
Immobilier & 2,5 (ventes 1,7 ; autres immeubles 0,4 ; bâti 0,4) ; 0,8 si la réaction des ventes n'est que temporaire [S]
& mobilité résidentielle accrue [E, étranger] ; aucun effet chiffré sur l'emploi ni sur la construction \\
Impôt sur les sociétés & 0,03 à 0,08 sur les équipements [S] & capital +11~Md€ (5,6 à 18,7) ; PIB potentiel
+0,06~\% (0,03 à 0,10), effet de niveau [S] \\
Plus-values & non chiffrée & moins de verrouillage aux âges élevés [E pour le sens] \\
\bottomrule
\end{tabularx}
\end{table}

\section{Le nouvel impôt sur le revenu}\label{ann:ircsg}
% Annexe du bloc IR-CSG-CRDS (note8p, 17/11/2026). Écrit par note8p/calculs/ircsg_annexe.py à partir de
% note8p/sorties/ircsg.json ; ne pas modifier à la main.

\subsection{Champ, recette et données}

\paragraph{Champ.} Le nouvel impôt remplace l'IR des revenus imposés au barème (salaires, pensions, chômage, bénéfices des indépendants, revenus fonciers), la CSG et la CRDS des revenus d'activité et de remplacement, et la CEHR due sur ces revenus. Les revenus du capital mobilier et les plus-values (PFU, option pour le barème comprise) relèvent du bloc capital. En 2025 : IR 86{,}6~Md€, CSG 135{,}8~Md€ (activité 108{,}5, pensions 26{,}7, chômage 0{,}6), CRDS 7{,}6~Md€, soit 229{,}9~Md€ ; s'y ajoutent la CEHR du barème (0{,}58~Md€ sur 1{,}39, au prorata des revenus du barème dans le revenu fiscal des tranches qui la paient, hypothèse) et les avantages maintenus pour l'invalidité et la dépendance (1{,}03~Md€, PLF 2026, Voies et moyens). Recette à tenir : 231{,}5~Md€.

\paragraph{Réconciliation de l'IR.} Comptes nationaux 2025 : 103{,}6~Md€ ; IRCOM des revenus 2024 porté en 2025 : 99{,}0 ; IR imputé par le modèle aux revenus du capital (12,8~\% de 90~Md€, convention de 2024) : 12{,}4 ; IR au PFU mesuré par le bloc capital (revenus 2025) : 8{,}4. L'écart de 4{,}0~Md€ n'est pas attribué : IR des revenus du capital imposés au barème par option (hors de notre périmètre), assurance-vie et plus-values immobilières comptées à tort dans les 90~Md€, erreur de répartition. S'il s'agissait en entier d'IR sur des revenus du barème, la recette à tenir monterait d'autant, soit $+$0{,}29~point sur tous les taux.

\paragraph{Modèle.} Foyers « purs » de \texttt{synthese/modele\_purs.py} : 25 tranches de l'IRCOM des revenus 2024, 60 foyers représentatifs par tranche, chacun scindé en foyers d'actifs et de retraités dans la proportion des pensions de la tranche ; couples homogènes, quatre partages des revenus (100/0, 88/12, 70/30, 55/45, Insee) ; 1,97~M de parents isolés (un enfant chacun) et deux enfants par couple « avec enfants » ; charge actuelle répartie entre types par des cas-types OpenFisca sans niche, puis recalée en 2025 composante par composante. Les cas-types sont recalculés avec OpenFisca-France 176.1 (droit des revenus 2025 ; IR, CEHR, CDHR, CSG, CRDS ; taux de CSG des pensions en régime stable) par \texttt{note8p/calculs/ircsg\_openfisca.py}.

\subsection{Taux marginaux d'aujourd'hui sur une assiette comparable}

Le taux marginal est le prélèvement supplémentaire sur 100~€ de revenu après cotisations en plus, calculé point par point par OpenFisca sur une grille de 100~€ (salarié seul du privé, sans enfant). Les taux nominaux ne s'additionnent pas : la CSG et la CRDS portent sur 98,25~\% du brut (100~\% au-delà de quatre plafonds), donc aussi sur les cotisations ; l'IR porte sur le revenu imposable, net de la CSG déductible (6,8 points) et des frais professionnels (10~\%, plafonnés). Au sommet, par euro de revenu après cotisations : entre 4 et 8 plafonds, IR 41{,}6 + CSG 10{,}3 + CRDS 0{,}56 = 52{,}4~\% ; au-delà de 8 plafonds, 51{,}7~\% ; plus 2{,}8 ou 3{,}7 points de CEHR. La CDHR ne joue pas pour un salarié (taux moyen supérieur à 20~\%). Médianes sur la grille : 54{,}7~\% entre 250\,000 et 450\,000~€, 55{,}7~\% au-delà de 550\,000~€ ; maximum 56{,}7~\%.

\begin{table}[h]\centering\footnotesize\setlength{\tabcolsep}{4pt}
\caption{Taux marginal sur 100~€ de revenu après cotisations, salarié seul (\%)}\label{tab:ircsg8_tm}
\begin{tabular}{rrrrrrr}\hline
Revenu (€) & Aujourd'hui & dont IR & dont CEHR & dont CSG-CRDS & Naïve & Proposition \\ \hline
10\,000 & 10{,}7 & 0{,}0 & 0{,}0 & 10{,}7 & 4{,}3 & 5 \\
15\,000 & 10{,}7 & 0{,}0 & 0{,}0 & 10{,}7 & 15{,}3 & 20 \\
19\,000 & 10{,}7 & 0{,}0 & 0{,}0 & 10{,}7 & 15{,}3 & 20 \\
20\,500 & 10{,}7 & 0{,}0 & 0{,}0 & 10{,}7 & 15{,}3 & 20 \\
25\,000 & 23{,}8 & 13{,}0 & 0{,}0 & 10{,}7 & 15{,}3 & 20 \\
30\,000 & 23{,}7 & 13{,}0 & 0{,}0 & 10{,}7 & 34{,}3 & 20 \\
38\,000 & 35{,}8 & 25{,}0 & 0{,}0 & 10{,}7 & 34{,}3 & 28 \\
45\,000 & 35{,}6 & 25{,}0 & 0{,}0 & 10{,}6 & 34{,}3 & 28 \\
60\,000 & 35{,}7 & 25{,}0 & 0{,}0 & 10{,}6 & 34{,}3 & 28 \\
90\,000 & 35{,}6 & 25{,}0 & 0{,}0 & 10{,}6 & 45{,}3 & 28 \\
110\,000 & 45{,}6 & 35{,}0 & 0{,}0 & 10{,}6 & 45{,}3 & 28 \\
150\,000 & 44{,}7 & 34{,}0 & 0{,}0 & 10{,}6 & 45{,}3 & 45 \\
200\,000 & 47{,}8 & 37{,}0 & 0{,}0 & 10{,}8 & 49{,}3 & 45 \\
250\,000 & 51{,}9 & 41{,}0 & 0{,}0 & 10{,}8 & 49{,}3 & 45 \\
300\,000 & 54{,}9 & 41{,}1 & 3{,}0 & 10{,}8 & 49{,}3 & 45 \\
400\,000 & 54{,}7 & 42{,}0 & 3{,}0 & 9{,}7 & 49{,}3 & 45 \\
600\,000 & 54{,}8 & 42{,}1 & 3{,}0 & 9{,}8 & 49{,}3 & 45 \\
\hline\end{tabular}
\par\footnotesize Fusion naïve : barème de l'IR (sans décote ni abattement) sur le revenu après cotisations, plus 4{,}3~\% sur tout le revenu (recette égale). Source : OpenFisca-France 176.1 ; calculs des auteurs.
\end{table}

\subsection{Choix du barème et variantes}

\paragraph{Exploration.} 1\,800 barèmes à quatre taux ont été essayés (\texttt{ircsg\_explore.py}, \texttt{ircsg\_affinage.py}), chacun calé sur la recette : taux d'entrée de 0 à 10~\% jusqu'à 6\,000 à 16\,000~€, deux taux intermédiaires, taux supérieur de 40 à 50~\% au-delà de 80\,000 à 150\,000~€. Le gain d'efficacité croît toujours quand le barème s'aplatit (le plus efficace à taux supérieur de 45~\%, 15~\% jusqu'à 100\,000~€ puis 45~\%, gagne 5,6~Md€ et rend 21~Md€ au dixième le plus aisé) : sans pondération d'équité, le modèle mène à un impôt proportionnel. Le choix est donc un choix de politique, guidé par cinq règles : un taux constant de 0,6 à 1,9 Smic (enseignement de Bozio et Wasmer, 2024) ; aucun salarié au Smic perdant ; taux supérieur de 45~\% sur le 1~\% le plus aisé ; chiffres ronds ; à progressivité égale, le barème le plus robuste à la perception de la CSG. Le seuil de 36\,500~€ est calé sur la recette (231{,}50~Md€, écart $-$0{,}03). Coût d'un point : tous les taux 14{,}0~Md€ ; 20 et 28~\% 8{,}3 ; 45~\% 0{,}42 ; 1\,000~€ de franchise par enfant 3{,}1~Md€.

\begin{table}[h]\centering\scriptsize\setlength{\tabcolsep}{3pt}
\caption{Variantes, recette constante (taux intermédiaires recalés)}\label{tab:ircsg8_var}
\begin{tabular}{lrrrrrrrrr}\hline
 & \multicolumn{3}{c}{Gain (Md€), CSG perçue} & Emplois & Retraités & Couple, 2 enf. & 1~\% & Mono-act. & Couple 1 rev. \\
 & entière & moitié & nulle & (milliers) & (Md€) & (€/foyer) & (Md€) & 60\,k€ (€) & 80\,k€ (€) \\ \hline
Proposition & 3{,}8 & 0{,}1 & $-$3{,}3 & 305 & $+$12{,}9 & $-$1\,130 & $-$2{,}2 & $+$1\,080 & $+$3\,603 \\
30~\% de 44 à 100~k€ & 3{,}4 & $-$0{,}3 & $-$3{,}8 & 306 & $+$12{,}3 & $-$1\,035 & $-$0{,}7 & $+$697 & $+$3\,548 \\
Seuil de 10\,000~€ à 0~\% & 3{,}1 & $-$0{,}6 & $-$4{,}1 & 408 & $+$12{,}4 & $-$1\,138 & $-$1{,}5 & $+$1\,272 & $+$4\,170 \\
Taux supérieur de 40~\% & 4{,}0 & 0{,}4 & $-$3{,}1 & 298 & $+$13{,}4 & $-$1\,239 & $-$3{,}9 & $+$1\,155 & $+$3\,755 \\
Taux supérieur de 50~\% & 3{,}6 & $-$0{,}1 & $-$3{,}6 & 311 & $+$12{,}4 & $-$1\,019 & $-$0{,}4 & $+$1\,007 & $+$3\,456 \\
Franchise de 4\,000~€/enfant & 3{,}9 & 0{,}2 & $-$3{,}2 & 284 & $+$11{,}0 & $-$635 & $-$2{,}2 & $+$1\,524 & $+$2\,927 \\
Franchise de 8\,000~€/enfant & 3{,}7 & 0{,}0 & $-$3{,}5 & 315 & $+$14{,}8 & $-$1\,608 & $-$2{,}3 & $+$429 & $+$4\,072 \\
Quotient conj. plafonné (2\,000~€) & 3{,}1 & $-$0{,}6 & $-$4{,}0 & 216 & $+$12{,}6 & $-$1\,330 & $-$2{,}0 & $-$690 & $+$2\,068 \\
Quotient conj. sans plafond & 3{,}1 & $-$0{,}6 & $-$4{,}0 & 174 & $+$12{,}9 & $-$1\,480 & $-$3{,}1 & $-$894 & $-$405 \\
Fusion naïve & 1{,}5 & $-$2{,}3 & $-$5{,}9 & 303 & $+$9{,}6 & $-$818 & $+$2{,}4 & $-$1\,427 & $-$969 \\
\hline\end{tabular}
\par Gain : variation de perte sèche réelle, jeu central ; perception sur les heures et l'effort, la décision de travailler se prenant sur le salaire net. Mono-act. : famille à un revenu, deux enfants, 60\,000~€ brut ; couple 1 rev. : sans enfant, 80\,000~€ brut.
\end{table}

\subsection{Résultats distributifs}

\paragraph{Déciles.} Variation moyenne par foyer fiscal (€ par an), déciles de foyers classés par revenu fiscal (IRCOM) : D1 $+$3 ; D2 $+$19 ; D3 $+$373 ; D4 $+$526 ; D5 $+$45 ; D6 $+$113 ; D7 $-$223 ; D8 $+$292 ; D9 $-$57 ; D10 $-$690. Le dixième le plus aisé paie 2{,}9~Md€ de moins, le 1~\% 2{,}2~Md€ de moins. Les déciles 3 et 4, où les retraités modestes sont nombreux, paient le plus. Actifs $-$11{,}3~Md€ ($-$416~€ par foyer) ; retraités $+$12{,}9~Md€ ($+$889~€) ; 33~\% des foyers paient plus de 500~€ de plus.

\paragraph{Couples et impôt individuel.} Variation moyenne selon la part du revenu du couple gagnée par le premier conjoint (100/0, 88/12, 70/30, 55/45) : couples d'actifs sans enfant $+$2\,945 / $+$1\,116 / $-$388 / $-$850~€ ; avec enfants $+$1\,166 / $-$492 / $-$1\,657 / $-$1\,745~€ ; couples de retraités $+$3\,620 / $+$2\,315 / $+$1\,116 / $+$767~€. Les perdants de l'individualisation sont les couples aux revenus très inégaux.

\begin{table}[h]\centering\footnotesize
\caption{Variation de l'impôt par configuration familiale, selon le revenu brut du ménage (€ par an)}\label{tab:ircsg8_fam}
\begin{tabular}{lrrrrrrr}\hline
 & 20\,k€ & 30\,k€ & 40\,k€ & 60\,k€ & 80\,k€ & 100\,k€ & 150\,k€ \\ \hline
Salarié seul & $-$160 & $-$130 & $-$480 & $-$1\,930 & $-$3\,290 & $-$4\,660 & $-$8\,660 \\
Couple biactif, sans enfant & $-$1\,020 & $-$930 & $-$320 & $-$560 & $-$500 & $-$2\,400 & $-$5\,900 \\
Couple à un revenu, sans enfant & $-$160 & $+$660 & $+$1\,480 & $+$2\,590 & $+$3\,600 & $+$2\,240 & $+$1\,210 \\
Couple biactif, 2 enfants & $-$1\,550 & $-$2\,130 & $-$2\,310 & $-$1\,080 & $-$1\,830 & $-$2\,150 & $-$5\,650 \\
Couple à un revenu, 2 enfants & $-$1\,620 & $-$1\,740 & $-$920 & $+$1\,120 & $+$1\,810 & $+$2\,490 & $-$570 \\
Couple biactif, 3 enfants & $-$1\,550 & $-$2\,330 & $-$2\,610 & $-$2\,280 & $-$1\,170 & $-$1\,880 & $-$3\,710 \\
Couple à un revenu, 3 enfants & $-$1\,910 & $-$2\,430 & $-$2\,120 & $-$460 & $+$1\,980 & $+$2\,710 & $+$1\,010 \\
Parent isolé, 2 enfants & $-$1\,620 & $-$1\,740 & $-$920 & $-$430 & $-$580 & $-$1\,950 & $-$7\,990 \\
Retraité seul & $+$1\,340 & $+$910 & $-$310 & $-$1\,910 & $-$3\,720 & $-$6\,080 & $-$10\,670 \\
Couple de retraités & $+$1\,000 & $+$1\,110 & $+$1\,240 & 0 & $-$1\,180 & $-$3\,350 & $-$7\,860 \\
\hline\end{tabular}
\par Cas-types OpenFisca (salariés du privé, couples mariés ou pacsés, biactifs à 60/40 ; retraités de plus de 65 ans, pension brute) ; charge d'aujourd'hui : IR, CEHR, CSG, CRDS.
\end{table}

\subsection{Efficacité et effets réels}

\paragraph{Méthode.} Celle de \texttt{revue/methode\_v2.md}, recopiée dans \texttt{ircsg\_modele.evalue} pour appliquer la perception de la CSG aux deux marges. Marge intensive : $\Delta Q = e\,g\,[\ln(1-t_1) - \ln(1-\tilde t_0)]$, où $t_1$ est le taux marginal du nouvel impôt (entièrement perçu), $\tilde t_0$ le taux actuel perçu (IR, plus la part perçue de la CSG-CRDS) ; chaque euro perdu coûte le coin moyen total non contributif (impôt, cotisations non contributives, retrait de la prime d'activité, inchangés). Élasticités du revenu réel 0,02/0,05/0,10 sous la médiane, 0,05/0,10/0,20 au-dessus (0,25 au sommet dans le jeu haut). Marge extensive : même formule sur le taux de prélèvement du premier salaire, adultes gagnant moins de deux Smic, $\eta$ = 0,10/0,25/0,50. Le revenu déclaré du dixième supérieur est un effet sur la recette (non compté).

\begin{table}[h]\centering\footnotesize
\caption{Gain d'efficacité (Md€ par an), emplois (milliers) et PIB potentiel, selon la perception de la CSG}\label{tab:ircsg8_eff}
\begin{tabular}{lrrrrrrr}\hline
CSG d'aujourd'hui perçue & \multicolumn{3}{c}{Gain : bas, central, haut} & \multicolumn{3}{c}{Emplois} & PIB (\%) \\ \hline
entière & 1{,}9 & 3{,}8 & 9{,}0 & 122 & 305 & 610 & $+$0{,}91 \\
moitié (heures, effort) & 0{,}0 & 0{,}1 & 0{,}8 & 122 & 305 & 610 & $+$0{,}32 \\
nulle (heures, effort) & $-$1{,}7 & $-$3{,}3 & $-$7{,}0 & 122 & 305 & 610 & $-$0{,}24 \\
moitié (deux marges) & $-$1{,}0 & $-$2{,}4 & $-$4{,}4 & $-$54 & $-$136 & $-$272 & $-$0{,}49 \\
nulle (deux marges) & $-$3{,}6 & $-$8{,}2 & $-$16{,}8 & $-$219 & $-$549 & $-$1\,097 & $-$1{,}80 \\
\hline\end{tabular}
\par PIB : Cobb-Douglas (part du travail 2/3), capital ajusté, jeu central ; à capital donné, deux tiers de ces chiffres.
\end{table}

Le gain central s'annule sous 48~\% de perception. Pensions futures : hausse du taux moyen sur les pensions de 3{,}1~points, soit un prélèvement implicite de 0{,}84~point sur le salaire ($\kappa = 0{,}274\,\Delta\tau$, part contributive des cotisations) ; gain 3{,}9 s'il n'est pas vu, 3{,}3 à moitié, 2{,}8 en entier.

\subsection{Retraités, trajectoire et limites}

\paragraph{Retraités.} Hausse de 12{,}9~Md€ (3{,}1~\% des pensions ; charge actuelle 50{,}2~Md€). Repères en droit actuel : abattement de 10~\% 5{,}32~Md€ ; abattement des plus de 65 ans 0{,}42 ; demi-part des anciens combattants 0{,}47 ; CSG de toutes les pensions au taux de 9,2~\% : $+$11{,}9~Md€ (approximation). Le nouveau barème rend une partie de ces sommes (taux de 5~\% jusqu'à 12\,000~€, 45~\% au sommet). Variation par retraité seul : P10 $+$27, médiane $+$845, P90 $+$1\,451~€ ; par couple : médiane $+$1\,479, P90 $+$3\,296~€.

\paragraph{Trajectoire.} Chaque foyer paie $(1-\lambda)$ fois l'ancien impôt et $\lambda$ fois le nouveau, $\lambda$ = 0,25 ; 0,5 ; 0,75 ; 1 de 2027 à 2030 : la recette est constante chaque année. En 2030 (Md€, niveau 2025) : retraités $+$12{,}88, actifs $-$11{,}30, CEHR du barème $-$0{,}58, invalidité et dépendance $-$1{,}03, total $-$0{,}05. L'extinction de la CRDS après 2033 (droit en vigueur) se traduit en 2034 par une baisse de 0{,}54~point de tous les taux, hors bilan de la réforme.

\paragraph{Limites.} Foyers synthétiques (pas de ménages mixtes actifs-retraités, pas de familles de trois enfants ou plus dans le modèle agrégé, niches réelles réparties au prorata) ; perception de la CSG non mesurée ; suppression des crédits d'impôt (emploi à domicile, garde d'enfants : 8,8~Md€) héritée du modèle ; IR non attribué de 4{,}0~Md€ ; IR de 2025 peut-être à un point haut ; répartition des revenus au PFU entre tranches inconnue ; pas de demande de travail ni de réaction des retraités ; frontaliers et non-résidents (CJCE, C-169/98) à expertiser.



\section{La fiscalité immobilière}\label{ann:foncier}
% Annexe du bloc « fiscalité immobilière » de la note de 8 pages (note8p/CADRE.md). Générée par
% note8p/calculs/foncier_annexe.py à partir de note8p/sorties/foncier.json (note8p/calculs/foncier_calculs.py, qui importe
% foncier/simulation.py). Compte rendu : note8p/sorties/foncier.md. Montants en Md€ de 2025, assiettes 2025 non indexées.

\subsection{Données, champ et calibrage}

Le modèle communal est celui de la note fusionnée (\texttt{foncier/simulation.py}, inchangé) : 30,5~millions de logements
des particuliers de métropole hors Alsace-Moselle ; valeur = prix médian au m² des ventes 2023-2024 de la commune × surface
médiane ; part du sol mesurée par les terrains à bâtir (maisons) ou par résidu (appartements), calée sur les comptes de
patrimoine des ménages (46,2~\% fin 2025) ; taxe foncière de 2025 (REI) répartie au prorata de la valeur dans chaque
commune. Les agrégats nationaux viennent des comptes de patrimoine de l'Insee (tableau 8.200, fin 2025) : biens privés
(sociétés non financières, sociétés financières, ménages, institutions sans but lucratif), hors terres agricoles et biens
publics. Sol : 6\,656~Md€, dont 565 de terrains non bâtis ; bâti (logements et bâtiments non résidentiels)
: 6\,560~Md€.

\textbf{Droits de mutation.} Taux permanent de 5,807~\% : 4,50~\% pour le département (art.~1594~D CGI), 1,20~\% pour la
commune (art.~1584) et 2,37~\% de la part départementale de frais d'assiette (art.~1647~V). La majoration à 5~\% (loi
n°~2025-127, art.~116), dont les primo-accédants sont exclus, cesse le 31~mars 2028 ; la référence de la réforme est donc
le taux permanent. France entière, au taux permanent et sur les ventes de 2025, les droits valent 13,1~Md€
sur les logements et 3,0~Md€ sur les autres biens (différence OFGL moins Compte du logement, ramenée au
taux permanent), soit 16,1~Md€ : la division par deux coûte \textbf{8,07~Md€}, contre
7,58 dans l'ancien calibrage, limité au champ du modèle. Le coût suit le cycle :
7,2~Md€ aux ventes moyennes de 2023-2024, 10,0~Md€ à celles de 2022. Un point de taux vaut
2,78~Md€.

\subsection{Architecture et formules}

Chaque bien $i$ a une base $b_i = S_i + B_i/4$ ($S$ : valeur du sol, bâti ou non ; $B$ : valeur du bâti). Il paie
$(t_c + t_n)\,b_i$ : le sol est taxé au taux $t$ et le bâti au taux $t/4$. Lors de la bascule, le taux communal
$t_c = \sum_{i \in c} \mathrm{TF}_i^0 / \sum_{i \in c} b_i$ garde la recette de chaque commune ; la commune le vote
ensuite librement, le rapport de 4 étant fixé par la loi. Le taux national uniforme $t_n$ lève l'enveloppe $E$ :
$t_n = E / \sum_i b_i$, la part des logements des ménages (64,5~\% de la base) étant répartie
entre les logements du modèle au prorata de $b_i$.

\textbf{Capitalisation.} La valeur du sol est la rente nette capitalisée : $S = \rho/(r + \tau_S)$, où $\tau_S$ est
l'impôt sur le sol rapporté à sa valeur. Aujourd'hui, $\tau_S^0 = 0{,}358~\%$ (taxe / valeur des logements, appliquée à
tous les biens bâtis) et 0,08~\% pour les terrains non bâtis (hypothèse). Après la réforme, $\tau_S^1 = t_c + t_n$ ; le
taux communal monte pour garder la recette sur un sol moins cher. Le point fixe
$S^1 = S^0 (r+\tau_S^0)/(r+\tau_S^1)$, $t_n = E/(S^1 + B/4)$ est résolu numériquement avec $r = 3~\%$ (variante 4~\%), le
bâti restant à son coût. La perte de valeur des biens est la valeur actuelle de la hausse, $\Delta \mathrm{TF}_i / r$
(baisse des droits gardée par l'acheteur, cas central ; variante symétrique dans le tableau~\ref{tab:fon8_env}).

\begin{table}[!htbp]\centering\scriptsize
\caption{Enveloppes de recette foncière nouvelle}\label{tab:fon8_env}
\begin{tabular}{@{}lrrrrrrrrr@{}}
\toprule
Enveloppe & \multicolumn{3}{c}{Taux national sur le sol (\%)} & Baisse du & Recette au & Payé par & \multicolumn{2}{c}{Perte de valeur (Md€)} & Net \\
\cmidrule(lr){2-4}\cmidrule(lr){8-9}
(Md€) & statique & $r$ = 3 \% & $r$ = 4 \% & sol (\%) & taux statique & logements & tous & logements & 2040 \\
\midrule
8,1 (substitution) & 0,097 & 0,106 & 0,103 & 9,8 & 7,4 & 5,2 & 269 & 173 & $-$0,8 \\
10 & 0,121 & 0,132 & 0,129 & 10,6 & 9,2 & 6,4 & 333 & 215 & +1,0 \\
15 (recommandée) & 0,181 & 0,201 & 0,196 & 12,4 & 13,5 & 9,7 & 500 & 322 & +5,5 \\
20 & 0,241 & 0,272 & 0,264 & 14,3 & 17,7 & 12,9 & 667 & 430 & +10,0 \\
\bottomrule
\end{tabular}
\par\vspace{2pt}{\raggedright\scriptsize Taux du bâti : un quart du taux du sol. « Recette au taux statique » : ce que
rapporterait le taux calculé sans capitalisation, une fois les valeurs ajustées. Perte de valeur à $r$ = 3~\% (4~\% : un
quart de moins). Si la baisse des droits se capitalise aussi, la perte nette de tous les propriétaires vaut
231~Md€ pour 15~Md€. Sans les terrains non bâtis, tous les taux montent de 7~\%.\par}
\end{table}

\subsection{Résultats pour les logements}

\begin{table}[!htbp]\centering\scriptsize
\caption{Taxe foncière des logements par dixième de prix au m² de la commune}\label{tab:fon8_dec}
\begin{tabular}{@{}rrrrrrrrrr@{}}
\toprule
 & Prix & Taxe/valeur & \multicolumn{4}{c}{Variation de taxe (€/an ; \% pour 15)} & Taxe $-$ droits & Taxe/valeur & Valeur \\
\cmidrule(lr){4-7}
Dixième & (€/m²) & 2025 (\%) & 10 & 15 & 15 (\%) & 20 & 15 (€/an) & 15 (\%) & 15 (\%) \\
\midrule
1 & 1\,035 & 0,54 & +36 & +60 & +13 & +84 & $-$21 & 0,61 & $-$2,4 \\
2 & 1\,404 & 0,49 & +61 & +97 & +17 & +132 & $-$9 & 0,57 & $-$2,7 \\
3 & 1\,693 & 0,49 & +79 & +124 & +18 & +169 & $-$6 & 0,58 & $-$2,9 \\
4 & 2\,000 & 0,46 & +125 & +184 & +24 & +243 & +38 & 0,57 & $-$3,7 \\
5 & 2\,356 & 0,43 & +161 & +235 & +28 & +310 & +70 & 0,56 & $-$4,1 \\
6 & 2\,737 & 0,41 & +185 & +274 & +31 & +363 & +87 & 0,54 & $-$4,3 \\
7 & 3\,176 & 0,39 & +233 & +340 & +37 & +448 & +131 & 0,54 & $-$4,8 \\
8 & 3\,749 & 0,36 & +278 & +413 & +42 & +548 & +182 & 0,51 & $-$5,0 \\
9 & 4\,746 & 0,31 & +331 & +496 & +49 & +662 & +202 & 0,46 & $-$5,0 \\
10 & 6\,261 & 0,18 & +513 & +780 & +93 & +1\,047 & +310 & 0,35 & $-$5,6 \\
\bottomrule
\end{tabular}
\par\vspace{2pt}{\raggedright\scriptsize Enveloppes de 10, 15 et 20~Md€ en régime permanent ; « taxe $-$ droits » :
hausse de taxe moins la baisse des droits rapportée à l'ensemble des logements de la commune ; valeur : perte de valeur
capitalisée ($r$ = 3~\%). Avec 15~Md€ : hausse médiane +237~€, 9\textsuperscript{e}
décile +655~€, 99\textsuperscript{e} centile +975~€ ;
40~\% des logements au-delà de +300~€, 12~\% au-delà de +600~€ ;
baisse pour 5~\% ; Paris +156~\% ; maisons
+263~€, appartements +351~€. Impôt sur le bâti des logements :
14,3~Md€ aujourd'hui, 7,3 avec la seule part communale à la nouvelle base, 9,4 avec 15~Md€.\par}
\end{table}

\subsection{Trajectoire 2027-2040}

Calendrier (choix) : loi en 2027, valeurs de référence au 1\textsuperscript{er}~janvier 2028, premier avis à l'automne
2029 ; part nationale et passage de la part communale à la nouvelle base par tiers en 2029, 2030 et 2031 ; taux des droits
abaissé au 1\textsuperscript{er}~janvier des mêmes années. Une clause suspend chaque palier de baisse tant que le palier
de taxe n'est pas mis en recouvrement. Hypothèses : 3~\% des ventes annuelles décalées de décembre à
janvier par point de baisse annoncée (les +26~\% de ventes du mois précédant la hausse de 0,7~point de 2014 représentent
environ 2~\% des ventes de l'année, soit 3~\% par point) ; report ouvert aux propriétaires occupants modestes, qui doivent
25~\% des hausses des occupants, utilisé par la moitié d'entre eux, avec un intérêt de 3~\% et un
remboursement de 1/15 du stock par an ; déductibilité l'année suivante au taux effectif de 20~\% pour
la part des sociétés soumises à l'IS (hors HLM, 18,7~\% de la base) et de
25~\% pour la hausse des bailleurs. Aucune reprise durable des ventes n'est comptée : à
6~\% de ventes par point, le coût des droits tomberait de 8,1 à 6,6~Md€.

\begin{table}[!htbp]\centering\scriptsize
\caption{Trajectoire budgétaire (Md€ ; + = recette publique)}\label{tab:fon8_traj}
\begin{tabular}{@{}lrrrrrrrrr@{}}
\toprule
 & 2027 & 2028 & 2029 & 2030 & 2031 & 2032 & 2033 & 2035 & 2040 \\
\midrule
Taxe foncière nouvelle & 0 & 0 & +5,0 & +10,0 & +15,0 & +15,0 & +15,0 & +15,0 & +15,0 \\
Droits de mutation & 0 & 0 & $-$2,7 & $-$5,4 & $-$8,1 & $-$8,1 & $-$8,1 & $-$8,1 & $-$8,1 \\
Achats différés & 0 & $-$0,5 & 0 & 0 & +0,2 & 0 & 0 & 0 & 0 \\
Report de paiement & 0 & 0 & $-$0,2 & $-$0,4 & $-$0,6 & $-$0,5 & $-$0,5 & $-$0,4 & $-$0,3 \\
Déductibilité & 0 & 0 & 0 & $-$0,4 & $-$0,8 & $-$1,2 & $-$1,2 & $-$1,2 & $-$1,2 \\
\textbf{Net, 15 Md€} & 0 & $-$0,5 & +2,1 & +3,8 & +5,8 & +5,2 & +5,2 & +5,3 & +5,5 \\
Net, 10 Md€ & 0 & $-$0,5 & +0,5 & +0,8 & +1,2 & +0,8 & +0,8 & +0,8 & +1,0 \\
Net, 20 Md€ & 0 & $-$0,5 & +3,7 & +6,9 & +10,3 & +9,6 & +9,7 & +9,8 & +10,0 \\
Net, substitution & 0 & $-$0,5 & $-$0,1 & $-$0,4 & $-$0,5 & $-$0,9 & $-$0,9 & $-$0,9 & $-$0,8 \\
Net, 15 Md€, baisse unique 2030 & 0 & 0 & +3,4 & +1,8 & +5,6 & +5,2 & +5,2 & +5,3 & +5,5 \\
Net, 15 Md€, un an de retard & 0 & 0 & $-$0,5 & +2,1 & +3,8 & +5,8 & +5,2 & +5,3 & +5,4 \\
Stock de report (15 Md€) & 0 & 0 & 0,2 & 0,6 & 1,2 & 1,8 & 2,3 & 3,4 & 5,6 \\
Taux des droits (\%) & 5,81 & 5,81 & 4,84 & 3,87 & 2,90 & 2,90 & 2,90 & 2,90 & 2,90 \\
\bottomrule
\end{tabular}
\par\vspace{2pt}{\raggedright\scriptsize Années 2034, 2036-2039 : voir \texttt{note8p/sorties/foncier.json}, table
\texttt{annuel} (15~Md€) et \texttt{annuel\_par\_scenario}. Le report est un prêt remboursable, compté ici en besoin de
financement net.\par}
\end{table}

\subsection{Cas-types}

\begin{table}[!htbp]\centering\scriptsize
\caption{Achats et taxe foncière (prix médian DVF au m² de la commune × surface)}\label{tab:fon8_cas}
\setlength{\tabcolsep}{3pt}
\begin{tabular}{@{}lrrrrrrrrrr@{}}
\toprule
 & Prix & Taux & Droits & Droits & Droits & Taxe & \multicolumn{2}{c}{Hausse (€/an)} & Bilan & Perte du \\
\cmidrule(lr){8-9}
 & (€) & 2026 (\%) & 2026 & permanents & après & actuelle & 15 & 10 & 7 ans & vendeur \\
\midrule
Orléans, app. 65 m² & 170\,000 & 5,81 & 9\,871 & 9\,871 & 4\,936 & 1\,102 & +298 & +223 & +3\,079 & $-$9\,934 \\
Lyon, app. 80 m² & 372\,000 & 6,32 & 23\,505 & 21\,601 & 10\,800 & 1\,242 & +665 & +450 & +6\,660 & $-$22\,151 \\
Massy, app. 65 m² & 269\,000 & 5,81 & 15\,620 & 15\,620 & 7\,810 & 925 & +475 & +326 & +4\,849 & $-$15\,844 \\
Ste-Geneviève, maison 100 m² & 365\,000 & 6,32 & 23\,063 & 21\,194 & 10\,597 & 1\,357 & +562 & +379 & +7\,098 & $-$18\,722 \\
\bottomrule
\end{tabular}
\par\vspace{2pt}{\raggedright\scriptsize Orléans et Massy : primo-accédants (hors majoration) ; Lyon et
Sainte-Geneviève-des-Bois : autres acheteurs, à 6,32~\% jusqu'au 31~mars 2028 (tableau DGFiP au 1\textsuperscript{er}~juin
2026). Bilan à 7 ans : économie de droits moins valeur actuelle de sept ans de hausse, prix inchangés. Perte du vendeur :
hausse capitalisée à 3~\%, qui revient à l'acheteur sous forme d'un prix plus bas si la taxe se capitalise. Mobilité : une
famille qui quitte Orléans pour un poste en Essonne paie 21\,194~€ de droits au taux permanent
(23\,063 en 2026), 10\,597~€ après réforme, soit l'équivalent de 294~€
puis 147~€ d'augmentation nette par mois pendant six ans.\par}
\end{table}

\subsection{Gains, sensibilités et variantes écartées}

\textbf{Perte sèche} (indicateur d'efficacité, triangle de Harberger, méthode de l'annexe de la note fusionnée). Ventes de
logements : $-$0,86, $-$1,71 et $-$2,29~Md€
pour 3, 6 et 8~\% de ventes en moins par point (soit +9, +18 et
+25~\% de ventes durables) ; zéro si la réaction n'est que temporaire. Autres immeubles :
$-$0,40~Md€ (transposition). Bâti des logements (élasticité 1, coût d'usage 4,5~\%) :
$-$0,41, $-$0,37 et $-$0,33~Md€ pour 10, 15 et 20~Md€. Le sol ne crée
pas de perte sèche. Schéma g6 : un terrain de 200\,000~€ et un immeuble de 600\,000~€ ; au taux unique de
0,358~\%, construire ajoute 2\,149~€ d'impôt par an ; avec
0,651~\% sur le sol et 0,163~\% sur le bâti (même recette), 977~€,
et le terrain nu paie 1\,302~€ au lieu de 716.

\textbf{Variantes écartées.} \emph{Substitution pure} (8,1~Md€) : le solde est négatif ($-$0,8~Md€ par an) à cause de la
déductibilité, et rien ne finance l'IS. \emph{Taux national sur le seul sol} : il évite les 20~\% de la part nationale
qui portent sur le bâti, mais le rapport sol/bâti varierait alors entre communes et il faudrait deux bases dans l'avis.
\emph{Baisse des droits en une fois} (2030) : aucune année négative dans le calcul, mais une incitation à différer les
achats pendant plus de deux ans que l'hypothèse d'un mois de décalage ne capte pas. \emph{Taux départemental} : de 0,13
à 0,79~\% du sol selon le département (ancien calcul). \emph{Suppression des droits} : 16,2~Md€ ; orientation de long
terme seulement. \emph{Baisse immédiate avant la taxe} (ancienne trajectoire) : 8,1~Md€ de dette par année d'avance.

\textbf{Expériences de référence.} \emph{Bade-Wurtemberg} (Landesgrundsteuergesetz du 4/11/2020, modifié en 2023 et
le 18/11/2025) : valeur = surface × \emph{Bodenrichtwert} de la zone, fixé par les comités d'experts (§~38 al.~1-2) ;
preuve d'une valeur inférieure de plus de 30~\% par expertise (§~38 al.~4) ; taux légal de 1,30~pour mille, réduit de 30~\% pour
les terrains majoritairement résidentiels, de 25~\% pour le logement social, de 10~\% pour les monuments (§~40) ;
\emph{Hebesatz} unique par commune (§~50 al.~4), taux distinct possible sur les terrains constructibles non bâtis
(§~50a) ; valeurs au 1/1/2022, impôt dû en 2025 ; modèle validé par le Bundesfinanzhof (II~R~26/24 et II~R~27/24,
communiqué du 20/05/2026), recours constitutionnels pendants (source associative). \emph{Pennsylvanie} : rapports
municipaux sol/bâti de 1,4 à 17,5 dans 18 villes en 2000, dont Pittsburgh (1913-2001), Scranton (1913), Harrisburg (1975),
Allentown (1997) (Banzhaf et Lavery, 2010, tableau~2) ; comté et district scolaire taxent au taux unique. Oates et Schwab
(1997) : permis réels +70~\% à Pittsburgh entre 1960-79 et 1980-89 contre $-$14~\% en moyenne dans 15 villes, attribués
surtout à la pénurie de bureaux (comparaison avant-après, un seul traité). Plassmann et Tideman (2000) : effet positif
sur les permis dans 15 villes à deux taux contre 204 témoins, sensible à la spécification (résumé et thèse lus). Banzhaf
et Lavery (2010) : différences premières avec effets fixes par secteur de recensement ; +5 à 6 points de croissance
décennale du nombre de pièces, surtout par plus de logements ; environ +3 points (pièces) et +2 points (logements) pour un
rapport de 4 (version de travail de 2008). Pittsburgh a abandonné le double taux en 2001 après une réévaluation qui avait
relevé la valeur du sol de 81~\% (Hughes, 2006).

\textbf{Limites.} Taux uniforme de la valeur dans chaque commune (la dispersion réelle multiplie les perdants) ; valeurs
de l'Insee sans géographie pour les sociétés ; contenu des « autres terrains » ; droits sur les autres biens par
différence de deux sources ; paramètres du report et de la déductibilité supposés ; capitalisation complète supposée
(contestée par une étude danoise non publiée) ; premier avis en 2029 ambitieux (le Bade-Wurtemberg a mis plus de quatre
ans en partant de valeurs de zone existantes).


\section{L'impôt sur les sociétés}\label{ann:is}
% Annexe technique du bloc IS (note8p, version 4 : taux de 25 %, CIR supprimé). Chiffres : note8p/sorties/is3.json
% (note8p/calculs/is3_reforme.py) et note8p/sorties/is4.json (note8p/calculs/is4_reforme.py) ; figure g7 :
% note8p/calculs/is3_figure_g7.py.
\subsection{Impôt sur les sociétés : données, méthode et trajectoire}\label{ann:is}

\paragraph{Droit de référence (2026).} Taux normal de 25~\%, maintenu ; contribution sociale de 3,3~\% au-delà de
763\,000~€ d'impôt. Taux réduit de 15~\% sur les 42\,500 premiers euros de bénéfice pour les sociétés dont le chiffre
d'affaires est d'au plus 10~M€ et dont le capital, entièrement libéré, est détenu à 75~\% au moins par des personnes
physiques (CGI, art.~219~I~b ; BOI-IS-LIQ-20-10 ; service-public.fr, fiche vérifiée le 17/02/2026). Contribution
exceptionnelle sur les grandes entreprises due au titre des deux premiers exercices clos à compter du 31/12/2025, puis
éteinte (LFI 2026, art.~12 ; BOI-IS-AUT-60) : droit en vigueur, non comptée. Amortissement dégressif : coefficient 2,25
au-delà de six ans (CGI, art.~39~A). Déficits : report en avant imputable à hauteur de 1~M€ plus 50~\% du bénéfice
au-delà (BOI-IS-DEF-10-30) ; report en arrière d'un exercice dans la limite de 1~M€ (BOI-IS-DEF-20-10). Charges
financières nettes déductibles jusqu'au plus élevé de 3~M€ et de 30~\% de l'EBITDA fiscal (art.~212~bis). CIR (art.~244
quater~B) : 30~\% des dépenses jusqu'à 100~M€, 5~\% au-delà ; depuis la LF 2025 (art.~55-56, d'après Légifiscal), forfait
de fonctionnement de 40~\% des dépenses de personnel, frais de brevets exclus ; CII ramené à 20~\% et prorogé jusqu'au
31/12/2027. Créance de CIR imputée sur l'IS de N+1 à N+3, reliquat restitué en N+4 ; remboursement immédiat pour les PME
au sens européen, les jeunes entreprises innovantes, les entreprises nouvelles et en difficulté (art.~199~ter~B, d'après
le Sénat, note de présentation « Remboursements et dégrèvements », PLF~2027).

\paragraph{Données et conventions.} Bénéfices imposables de l'exercice 2025 (DGFiP, \emph{Statistiques} n°~51) : 256,8~Md€
au taux normal, 30,6~Md€ au taux réduit, soit 267,2 et 31,8~Md€ en 2027 ; +3~\% par an ensuite. Un point de taux normal :
2,99~Md€ en 2027. Investissement éligible (machines, équipements, informatique des sociétés non financières, Insee
8.212) : 84,1~Md€ en 2025, 89,1~Md€ en 2027 ; stocks de fin 2025 : 541~Md€ (équipements), 400~Md€ (bâtiments).
Hypothèses : 75~\% de l'investissement éligible fait par des sociétés imposables ; plan actuel moitié linéaire, moitié
dégressif sur 7~ans. Années d'exercice ; en caisse, un an plus tard. Dépenses fiscales : Voies et moyens, tome~II,
PLF~2026 (coûts 2024 réalisés, 2025 et 2026 prévus), portées à +3~\% par an.

\paragraph{Formules.} Nouveau plan d'amortissement : $b\,e_0 + (1-b)\,d$, où $d_k$ est le plan ordinaire et $b$ la
fraction déduite l'année de l'achat. Coût de l'année $t$ : $C_t = \tau\, s \sum_{v\le t} I_v\,[\,b_v e_0 + (1-b_v) d - d\,]_{t-v}$ ;
en régime établi, $C_t = \tau\, s\, b\, I_t\,(1 - z(g))$, avec $z(r)=\sum_k d_k/(1+r)^k$ ; coût actualisé d'une année
d'achats : $\tau\, s\, b\, I\,(1-z(r))$. Coût d'usage relatif : $(1-\tau z)/(1-\tau)$. TMEI (Devereux et Griffith) :
$\tilde p = (1-\tau z)(\rho-\pi+\delta(1+\pi))/((1-\tau)(1+\pi)) - \delta$, $\mathrm{TMEI} = (\tilde p - r)/\tilde p$ ; $r = 5~\%$,
$\pi = 2~\%$, $\delta$ = 14~\% (équipements, 7~ans), 33~\% (logiciels, 3~ans), 3,1~\% (bâtiments, 25~ans) ; $\rho = 7{,}1~\%$
(fonds propres) ou $i(1-\tau)$ (emprunt) ; mixte : un tiers d'emprunt. Dépense de R\&D passée en charges avec un crédit
non imposable au taux $c$ : $\mathrm{TMEI} = -c/(1-\tau-c)$. Effets réels : $\Delta K_j/K_j = -\varepsilon_j\, s\, \Delta\ln c_j$,
$\varepsilon$ = 1 (0,5 à 1,6) pour les équipements, 0,5 (0,25 à 1, hypothèse) pour les bâtiments ; Cobb-Douglas à emploi
donné, $\Delta Y/Y = \sum_j \sigma_j \Delta K_j/K_j$, $\sigma_j = (r+\delta_j)K_j/Y$ (3,4~\% et 1,1~\%) ; investissement de
régime établi : $(\delta_j + g)\,\Delta K_j$.

\paragraph{Taux réduit.} Rendement à plein en 2027 : 3,18~Md€ statiques (10~points sur 31,8~Md€) ; 2,15~Md€ en série
prudente (coût officiel de 2019, 2,2~Md€ pour 13~points, ramené à 10~points et actualisé), retenue pour le bouclage.
Calendrier : 20~\% pour l'exercice 2027 (moitié du rendement), 25~\% dès 2028. Perte par société : $0{,}10\min(B,
42\,500)$, au plus 4\,250~€ ; 760\,000 bénéficiaires en 2019, au moins 720\,000 sociétés en équivalent plafond en 2025.

\begin{table}[!htbp]\centering\scriptsize
\caption{Trajectoire pour le bouclage (Md€ courants, années d'encaissement ; + = recette publique)}\label{tab:is8_annuel}
\setlength{\tabcolsep}{3.2pt}
\begin{tabular}{@{}lrrrrrrrr@{}}
\toprule
& Fin du & Déduction & \multicolumn{2}{c}{Niches hors CIR} & \multicolumn{2}{c}{Suppression du CIR} & \emph{Brevets} & Solde \\
\cmidrule(lr){4-5}\cmidrule(lr){6-7}
Année & taux réduit & de 50~\% & statique & prudent & statique & prudent & \emph{(hors)} & prudent \\
\midrule
2028 & 1,08 & -- & -- & -- & -- & -- & -- & 1,08 \\
2029 & 2,22 & -- & 1,77 & 1,47 & 7,66 & 4,22 & 1,33 & 7,91 \\
2030 & 2,29 & $-$3,92 & 2,30 & 1,88 & 7,89 & 4,64 & 1,38 & 4,89 \\
2031 & 2,35 & $-$7,16 & 2,46 & 2,01 & 8,13 & 4,98 & 1,42 & 2,18 \\
2032 & 2,42 & $-$5,71 & 2,63 & 2,15 & 8,37 & 6,70 & 1,46 & 5,56 \\
2033 & 2,50 & $-$4,54 & 2,81 & 2,29 & 8,62 & 6,90 & 1,50 & 7,14 \\
2034 & 2,57 & $-$3,56 & 2,99 & 2,44 & 8,88 & 7,10 & 1,55 & 8,55 \\
2035 & 2,65 & $-$2,63 & 3,18 & 2,59 & 9,15 & 7,32 & 1,59 & 9,93 \\
2036 & 2,73 & $-$1,67 & 3,39 & 2,76 & 9,42 & 7,54 & 1,64 & 11,35 \\
2037 & 2,81 & $-$1,04 & 3,49 & 2,84 & 9,70 & 7,76 & 1,69 & 12,38 \\
2038 & 2,89 & $-$0,91 & 3,59 & 2,93 & 9,99 & 8,00 & 1,74 & 12,91 \\
2039 & 2,98 & $-$0,93 & 3,70 & 3,01 & 10,29 & 8,23 & 1,79 & 13,29 \\
2040 & 3,07 & $-$0,96 & 3,81 & 3,10 & 10,60 & 8,48 & 1,85 & 13,69 \\
\bottomrule
\end{tabular}
\par\vspace{3pt}\parbox{\linewidth}{\raggedright\scriptsize Encaissement : exercice + 1, sauf CIR (profil des créances).
Taux réduit : série prudente (20~\% pour l'exercice 2027, 25~\% ensuite). Déduction : 25~\% pour l'exercice 2029, 50~\%
ensuite, coût brut. Niches hors CIR : mécénat, crédits culture, zonages, crédit famille, dès l'exercice 2028 ; prudent :
après réaction (facteurs 0,85, 0,75, 0,80, 0,90). CIR supprimé pour les dépenses de 2028 et suivantes ; statique : coût
total de l'année ; prudent : millésimes de créances 2028 et suivants seulement (le stock antérieur est imputé ou restitué
jusqu'en 2031), décote de 20~\%. Solde prudent : colonnes 2, 3, 5 et 7. Brevets : statique, non comptés. Mémoire, non
comptés (droit en vigueur) : fin de la contribution exceptionnelle, environ $-$7~Md€ par an par rapport à 2026 ;
extinction du CII. Source : \texttt{note8p/calculs/is4\_reforme.py}.}
\end{table}

\paragraph{Déduction immédiate.} Coût permanent : 0,67~Md€ de 2027 par an (0,54 à 0,81 selon le plan actuel), 0,99~Md€ en
2040 ; coût actualisé d'une année d'achats : 1,07~Md€ à 5~\% (0,73 au taux de l'État de 3,25~\%). Le coût des premières
années est une avance d'impôt, récupérée ensuite. Véhicules utilitaires (moitié du matériel de transport, hypothèse) :
coûts multipliés par 1,25. Aucune recette de croissance n'est comptée.

\paragraph{Niches : paquet chiffré.} Coûts statiques des Voies et moyens (2026). Mécénat des entreprises (210309,
1\,730~M€, 157\,399~bénéficiaires) : réduction remplacée par la déduction du don, $1{,}73\,(1-0{,}25/0{,}60) = 1{,}01$~Md€ à
dons inchangés. Crédits culture (320121, 129, 140, 135, 142, 128, 145, 148) : 0,71~Md€, moitié en 2028. Zonages (220104,
220109, 220102, 230602, 230606, 230608, 230605, 210305) : 0,66~Md€, fermés aux entrants, extinction linéaire sur huit ans
(hypothèse). Crédit famille (210308) : 0,22~Md€. Noyau hors CIR : 2,60~Md€ (2026). Rendement réalisable prudent
(hypothèses) : mécénat $\times 0{,}85$ (taux de 40~\% au-delà de 2~M€ de dons et report des excédents sur cinq ans ; la
baisse des dons accroît au contraire la recette) ; culture $\times 0{,}75$ (tournages perdus, recettes de TVA et de
cotisations en moins) ; zonages $\times 0{,}80$ (substitution) ; crédit famille $\times 0{,}90$. Hors paquet : taxe au
tonnage (régime européen, base mobile), outre-mer (à budgétiser), PTZ et HLM (logement), SIIC et SPPICAV (transparence).
La quote-part de frais et charges et le régime mère-fille éliminent la double imposition : ils ne sont pas des niches.
Doubles comptes : les dispositifs communs IR-IS sont comptés ici en entier (le bloc IR ne les reprend pas) ; les
exonérations de fiscalité locale en zone et le mécénat des ménages relèvent d'autres blocs.

\paragraph{Suppression du CIR.} Dépenses de R\&D engagées à compter de 2028 : plus de crédit, déduction ordinaire des
charges ; équipements de recherche dans la déduction de 50~\% (déjà dans l'investissement éligible). L'ancien recentrage
(plafond au niveau du groupe et fin du taux de 5~\%, 0,97~Md€) est absorbé : pas de double compte. Coût : 7,3~Md€
exécutés en 2025, plus de 7,4~Md€ projetés en 2027 (Sénat, PLF~2027) ; 8,0~Md€ en 2026 selon les Voies et moyens du
PLF~2026 (borne haute). Créance 2023 : 7,85~Md€, dont recherche 7,32 (PME 2,24, ETI 2,10, grandes entreprises 2,98) ;
33~bénéficiaires dépassant 100~M€ de dépenses au niveau de l'intégration fiscale reçoivent 25~\% de la créance recherche
(MESRE, novembre 2025). En 2024, sur 7,7~Md€ (CIR et CII), 4,7 restitués et 3,0 imputés (Sénat, d'après la DGFiP). Profil
retenu d'un millésime de créances (hypothèse construite sur ces chiffres) : 67~\% en N+1 (remboursements immédiats,
environ 36~\%, et 80~\% des imputations), 5~\% en N+2, 3~\% en N+3, 25~\% restitués en N+4. Réalisable : millésimes 2028 et
suivants seulement, $\times 0{,}8$ (emploi de recherche perdu, soutiens recréés hors de l'impôt ; hypothèse, aucune
élasticité ne la calibre). Effet sur la R\&D : un euro de CIR produit 0,9 à 1,5~euro de R\&D dans les PME (CNEPI, cité
par le Sénat ; Bozio, Cottet et Py, 2019), aucun effet établi pour les ETI et grandes entreprises (CNEPI 2021, IPP 2021).
R\&D des PME : $-$2,1 à $-$3,4~Md€ par an ; stock $-$11 à $-$19~Md€ (dépréciation de 15~\%, croissance de 3~\%) ; PIB
potentiel $-(r+\delta)\Delta S/Y$ = $-$0,07 à $-$0,12~\%, double si le rendement social est double du rendement privé
(Bloom, Schankerman et Van Reenen, 2013). Non compté. Variante : crédit conservé pour les seules PME, recette statique
5,1~Md€ (2027). TMEI de la R\&D : $-$67~\% avec le CIR à 30~\%, nul sans.

\begin{table}[!htbp]\centering\scriptsize
\caption{Taux marginaux effectifs d'un investissement neuf (\%)}\label{tab:is8_tef}
\setlength{\tabcolsep}{4pt}
\begin{tabular}{@{}lrrrrrr@{}}
\toprule
& \multicolumn{3}{c}{Équipement} & Logiciel & \multicolumn{2}{c}{Bâtiment} \\
\cmidrule(lr){2-4}\cmidrule(l){6-7}
Configuration & fonds propres & emprunt & 1/3 dette & fonds propres & fonds propres & emprunt \\
\midrule
Aujourd'hui (25~\%) & 18,0 & $-$23,8 & 7,6 & 19,5 & 22,0 & $-$19,3 \\
Taux à 22,5~\% (comparaison) & 16,0 & $-$20,8 & 6,5 & 17,5 & 19,7 & $-$16,7 \\
Déduction de 50~\% (proposition) & 9,9 & $-$37,0 & $-$1,7 & 10,8 & 22,0 & $-$19,3 \\
Déduction intégrale (long terme) & 0,0 & $-$53,4 & $-$13,1 & 0,0 & 22,0 & $-$19,3 \\
Petite société sous 42\,500~€, aujourd'hui & 10,4 & $-$12,8 & 3,8 & 11,4 & 13,0 & $-$10,0 \\
\bottomrule
\end{tabular}
\par\vspace{3pt}\parbox{\linewidth}{\raggedright\scriptsize Bâtiments hors champ de la déduction. Dépense de R\&D passée
en charges : TMEI nul ; avec le CIR à 30~\% : $-$67~\% ; au taux de 5~\% au-delà de 100~M€ : $-$7~\%. Avantage de la dette (écart
de coût du capital entre fonds propres et emprunt) : 2,06 points aujourd'hui, 1,90 avec la déduction de 50~\%, 1,74 avec la
déduction intégrale. Niveaux non comparables à ceux de l'OCDE ; seuls les écarts comptent.}
\end{table}

\paragraph{Dette.} Coût du capital d'une machine (rendement exigé 5~\%) : fonds propres 6,10, 5,55 et 5,00~\% ; emprunt
4,04, 3,65 et 3,26~\% pour une fraction déduite tout de suite de 0, 50 et 100~\%. L'écart diminue (une déduction anticipée
vaut plus actualisée à 7,1~\% qu'à $i(1-\tau) = 5{,}3$~\%), mais le coût de l'emprunt passe sous le rendement exigé : les
intérêts nominaux restent déduits sur tout le prix alors que l'État en a déjà pris $\tau b$ à sa charge. Corrections
(TMEI pour $b = 0{,}5$) : rabot général de 21~\% des charges financières nettes pour ramener le TMEI de l'emprunt à
$-$24~\%, de 14~\% pour un TMEI mixte nul ; déduction des seuls intérêts réels ($i-\pi$) : $-$19,5~\% (emprunt), +1,8~\% (mixte) ;
au-delà du plafond de l'art.~212~bis, l'intérêt marginal n'est pas déductible et le TMEI de l'emprunt égale celui des
fonds propres (9,9~\%). Le seuil de 3~M€ de charges nettes correspond à 50 à 75~M€ de dette (taux de 6 à 4~\%), celui de
1~M€ (Espagne) à 17 à 25~M€ : l'abaisser ne touche pas l'investissement marginal des PME. Recommandation : pas de
garde-fou dans la proposition centrale ; neutralité à traiter avec la déduction intégrale.

\paragraph{Pertes.} Valeur actualisée de la déduction de 50~\% pour 100 d'achat : 1,60 pour une société imposable ;
1,07, 0,41 et 0,09 si elle reste déficitaire 1, 3 ou 5~ans (déficits non indexés, imputés sans plafond) ; nulle au-delà
de la durée du plan. Part de l'investissement éligible faite par des sociétés déficitaires : non mesurée (aucune source
croisant investissement et situation fiscale) ; repère : 113~Md€ de déficits déclarés en 2022 pour 278~Md€ de bénéfices
(DGFiP n°~51) ; l'hypothèse de 75~\% d'investissement imposable porte coûts et effets.

\paragraph{Effets réels.} Coût d'usage des équipements : $-$2,1~\% avec la déduction de 50~\%, $-$4,2~\% avec la déduction
intégrale, $-$0,5~\% avec un taux de 22,5~\% (qui baisse aussi de 1,6~\% celui des bâtiments). Capital supplémentaire
(central, bornes) : 8,4 (4,2 à 13,4)~Md€ ; 17,0 (8,5 à 27,1) pour la déduction intégrale ; 4,6 (2,3 à 8,3) pour le taux à
22,5~\%. Investissement de régime établi : +1,4~Md€ par an. PIB potentiel, effet de niveau : +0,05~\% (0,03 à 0,09), soit
1,6~Md€, pour la déduction de 50~\% ; +0,11~\% pour la déduction intégrale ; +0,02~\% pour le taux à 22,5~\%. Par Md€ de coût
économique annuel (coût actualisé d'une année d'achats, ou coût permanent du taux), 7,8~Md€ de capital pour la
déduction, 0,6 pour le taux : calcul partiel, qui ignore la localisation et les bénéfices transférés ; Chodorow-Reich,
Zidar et Zwick (2024) trouvent un écart plus faible, de même sens. Petites sociétés sous le plafond : coût d'usage des
équipements $-$0,1~\%, des bâtiments +6~\% ; non agrégé, faute de connaître leur part du capital ; effet de trésorerie non
chiffré (sensibilité non estimée en France). Ces effets ne mesurent que la déduction : la fin
du taux réduit (bâtiments et trésorerie des petites sociétés) et la suppression du CIR (R\&D des PME) jouent en sens
inverse et ne sont pas chiffrées dans le +0,05~\% ; les autres niches n'ont qu'un effet de localisation ou de secteur.

\paragraph{Long terme, non retenu.} Passer de 50 à 100~\% en 2035 : 9,4~Md€ de plus la première année, 7,5 la deuxième,
puis 0,67~Md€ par an (euros de 2027) ; d'un coup en 2027 : 14,8~Md€ ; TMEI mixte de $-$13~\%, d'où un traitement des intérêts
à prévoir. Déduction pour fonds propres et remboursement des pertes : non chiffrés (fonds propres fiscaux non disponibles
dans une source vérifiée).

\paragraph{Cas types.} Machine de 500\,000~€ achetée le 1\textsuperscript{er} juillet, 7~ans, à 25~\% : déduction de la
première année de 35\,714 à 267\,857~€ (linéaire), soit 58\,036~€ de trésorerie, rendus ensuite ; avantage actualisé de
9\,544~€ (5~\%) ou 14\,155~€ (8~\%) ; dégressif : 52\,455 et 6\,461~€. Suramortissement « industrie~4.0 » (PLF~2027, art.~9,
proposé) sur la même machine : 43\,600~€ (40~\%), 87\,200~€ (80~\%), 130\,800~€ (120~\%) actualisés, contre 8\,003~€ pour la
déduction de 50~\% (plan moyen). Société de services : perte de $0{,}10\min(B, 42\,500)$, soit 3\,000~€ à 30\,000~€ de
bénéfice ; la déduction sur 10\,000~€ d'informatique annuelle lui vaut environ 90~€.

\paragraph{Variantes écartées et hypothèses retirées.} Taux à 22,5~\% : 7,5~Md€ par an pour un effet sur le capital
domestique moindre. Recentrage du CIR (plafond au niveau du groupe, fin du taux de 5~\%), remplacé par la suppression. Retirés par parcimonie : rendement « central » du taux
réduit (fuite de 16~\%), effet de trésorerie des PME, multiplicateur d'investissement de court terme, retour par les
bénéfices déclarés, sensibilités de la part imposable ; aucune décote comportementale sur les niches (la boîte à brevets
est exclue du bouclage prudent).

\paragraph{Limites.} Part imposable de l'investissement et plan actuel (non mesurés) ; élasticité des bâtiments ;
profil d'encaissement du CIR et décotes de réaction (hypothèses) ; charges financières nettes déduites non chiffrées ; tome~II des Voies et moyens du PLF~2027 non lu ; durées d'extinction des
zonages et taux du mécénat (CGI, art.~238~bis) non revérifiés sur les textes ; issue du PLF~2027 inconnue au 18/11/2026.


\section{La fiscalité du capital}\label{ann:pfu}
% Annexe du bloc « fiscalité du capital », révision v3 (note8p/CONSIGNES_V3.md). Chiffres : note8p/sorties/pfu3.json et
% pfu3_calculs.txt (programme note8p/calculs/pfu3_neutralite.py) ; coût d'un point de PFU, CEHR et CDHR :
% note8p/sorties/pfu.json (programme note8p/calculs/pfu_transmission.py). Figure : note8p/figures/g9_exemple_pv.

\subsection{La règle : reprise du prix, correction à la marge}

Le PFU reste à 31,4~\% : 12,8~\% d'IR (CGI, art.~200~A) et 18,6~\% de prélèvements sociaux (LFSS pour 2026,
art.~12). À la donation et au décès, la plus-value des titres de l'art.~150-0~A n'est plus effacée : le bénéficiaire
reprend le prix d'acquisition historique et n'est imposé que s'il vend [C]. Reste à éviter que l'impôt latent paie
aussi des droits de mutation. Notons $V$ la valeur transmise, $G$ la plus-value, $t$ le PFU et $D(\cdot)$ les droits
du bénéficiaire sur sa part. Le parent qui vend paie $tG + D(V-tG)$ ; l'héritier qui reçoit puis vend paie $D(V) +
t(G-a)$, où $a$ est la majoration du prix repris. Les deux charges sont égales si et seulement si
\[
t\,a \;=\; D(V) - D(V - tG),
\]
c'est-à-dire si l'héritier récupère, à la vente, les droits qu'il a payés sur l'impôt latent. Les droits ne doivent
pas frapper la part de l'héritage qui reviendra au fisc. D'où la règle retenue [C], dite « correction à la marge » :
\emph{le prix repris est majoré de $a = [D(V) - D(V - tG)]/t$}. Le notaire calcule deux fois les droits du bénéficiaire,
une fois sur sa part, une fois sur sa part diminuée de l'impôt latent. Il inscrit la différence dans la déclaration,
qui doit de toute façon porter le prix historique. La majoration s'attache aux titres : elle suit les ventes
partielles au prorata et se cumule si les titres sont transmis à nouveau. Elle équivaut à un crédit d'impôt à la vente,
mais sous la forme d'un prix, plus simple à suivre.

\begin{table}[!htbp]\centering\scriptsize
\caption{Règles comparées}\label{tab:cap_regles}
\begin{tabular}{@{}p{4.4cm}p{3.1cm}p{3.2cm}p{3.2cm}@{}}
\toprule
Règle & Neutralité vendre / transmettre & Si l'héritier ne vend jamais & Coût et simplicité \\
\midrule
(i) Taux moyen des droits $\times G$ ajouté au prix (v2 ; 26 U.S.C. §~1015(d)(6)) & Double imposition résiduelle
de $tG(m-\bar\tau)$ : 6~\% de la charge à 1~M€ & Droits inchangés & Simple ; recette de référence \\
(ii) Même correction au taux marginal $m$ & Exacte tant que $V$ et $V-tG$ sont dans la même tranche & Droits inchangés
& Simple ; 0,2 à 0,3~Md€/an de moins que (i) \\
(iii) Impôt latent déduit de l'assiette des droits & Exacte & Droits \emph{réduits} pour un impôt jamais payé &
Coûte 0,8 à 1,1~Md€/an dès 2027 \\
\textbf{(iv) Crédit à la vente des droits afférents à l'impôt latent (retenue, sous forme de prix)} & Exacte & Droits
inchangés & Un second calcul de droits ; comme (ii) \\
(v) Prix majoré de 30~\% de $G$ & Juste si $m = 30~\%$ ; sinon jusqu'à $-30~\%$ (conjoint) ou $+10~\%$ (neveu) &
Droits inchangés & Simple mais arbitraire \\
(vi) Réalisation au décès, impôt déduit des droits (Canada) & Exacte & Impôt dû quand même & Contraire à
l'imposition à la vente seule (non retenue) \\
\bottomrule
\end{tabular}
\end{table}

\begin{table}[!htbp]\centering\scriptsize\setlength{\tabcolsep}{4pt}
\caption{Charge totale (plus-value + droits) et écart avec « vendre puis transmettre »}\label{tab:cap_cas}
\begin{tabular}{@{}lrrrrrr@{}}
\toprule
 & Vendre, puis & Transmettre & \multicolumn{4}{c}{Réforme : écart (transmettre puis vendre)} \\
\cmidrule(l){4-7}
Cas (droit de 2026) & transmettre & aujourd'hui & (i) moyen & (iv) retenue & (iii) & (v) 30~\% \\
\midrule
100~k€ devenus 1~M€, un enfant ($\bar\tau$ 21~\%, $m$ 30~\%) & 410\,782 & 212\,962 & +24\,597 & 0 & 0 & 0 \\
250~k€ devenus 400~k€, un enfant (15~\%, 20~\%) & 95\,874 & 58\,194 & +2\,568 & 0 & 0 & $-$4\,710 \\
2~M€ devenus 40~M€, deux enfants (44~\%, 45~\%) & 23\,997\,389 & 17\,434\,789 & +168\,603 & 0 & 0 & +1\,789\,800 \\
100~k€ devenus 1~M€, un frère (44~\%, 45~\%) & 595\,818 & 440\,388 & +2\,716 & 0 & 0 & +42\,390 \\
100~k€ devenus 1~M€, un neveu (55~\%, 55~\%) & 672\,788 & 545\,618 & +1\,238 & 0 & 0 & +70\,650 \\
100~k€ devenus 1~M€, conjoint (exonéré) & 282\,600 & 0 & 0 & 0 & 0 & $-$84\,780 \\
\bottomrule
\end{tabular}
\par\vspace{3pt}\parbox{0.95\linewidth}{\raggedright\scriptsize Lecture : aujourd'hui, transmettre avant de vendre
divise la charge par deux dans l'exemple ; avec la règle (iv), les deux chemins coûtent autant. L'écart de la règle (i)
vient de la progressivité et de l'abattement : il disparaît presque aux taux plats (neveu) ou au taux maximal
(très grande fortune). Si l'héritier ne vend jamais, il paie avec (i), (ii), (iv) et (v) les droits d'aujourd'hui ;
avec (iii), moins qu'aujourd'hui (128\,182~€ au lieu de 212\,962~€ dans l'exemple).}
\end{table}

\subsection{Recettes 2027-2040 de la règle retenue}

Méthode, jeux d'hypothèses et calendrier sont ceux de la version 2 (\texttt{pfu\_transmission.py}) : impôt latent
transmis par an $L = 31{,}4~\%\,E\,b\,\lambda + 30~\%\,S/T$ (3,65 / 7,49 / 14,68~Md€ dans les jeux bas, central,
haut), part $a$ vendue dans les douze mois puis $h$ par an, impôt encaissé l'année qui suit la vente. Seul change le
facteur de correction : $(1-m)$ au lieu de $(1-\bar\tau)$, avec $\bar\tau$ = 15 / 10 / 5~\% [H] et
$m = 1{,}4\,\bar\tau$ [H]. Le rapport 1,4 est celui des cas-types en ligne directe (1,0 au taux maximal, 1,4 à 1~M€, 1,6
à 1,8 sous pacte Dutreil). Unités : Md€ au niveau du PIB de 2025.

\begin{table}[!htbp]\centering\scriptsize\setlength{\tabcolsep}{3pt}
\caption{Recette encaissée par année (Md€)}\label{tab:cap_annuel}
\begin{tabular}{@{}l*{14}{r}r@{}}
\toprule
 & 27 & 28 & 29 & 30 & 31 & 32 & 33 & 34 & 35 & 36 & 37 & 38 & 39 & 40 & Cumul \\
\midrule
\multicolumn{16}{@{}l}{\emph{Règle retenue (iv)}} \\
Bas & 0 & 0,19 & 0,44 & 0,49 & 0,54 & 0,58 & 0,63 & 0,67 & 0,72 & 0,76 & 0,80 & 0,85 & 0,89 & 0,93 & 8,5 \\
Central & 0 & 0,82 & 1,83 & 2,01 & 2,19 & 2,36 & 2,52 & 2,68 & 2,83 & 2,97 & 3,11 & 3,24 & 3,37 & 3,50 & 33,4 \\
Haut & 0 & 2,51 & 5,53 & 6,02 & 6,48 & 6,91 & 7,31 & 7,69 & 8,05 & 8,38 & 8,70 & 9,00 & 9,28 & 9,54 & 95,4 \\
\textbf{Prudent} (bas, deux ans plus tard) & 0 & 0 & 0 & 0,19 & 0,44 & 0,49 & 0,54 & 0,58 & 0,63 & 0,67 & 0,72 & 0,76 & 0,80 & 0,85 & 6,7 \\
\multicolumn{16}{@{}l}{\emph{Pour comparaison, jeu central}} \\
(i) taux moyen (v2) & 0 & 0,86 & 1,91 & 2,10 & 2,29 & 2,47 & 2,64 & 2,80 & 2,96 & 3,11 & 3,26 & 3,40 & 3,53 & 3,66 & 35,0 \\
(iii) déduction des droits & $-$1,05 & $-$0,10 & 1,08 & 1,29 & 1,50 & 1,69 & 1,88 & 2,07 & 2,24 & 2,41 & 2,57 & 2,72 & 2,87 & 3,02 & 24,2 \\
\bottomrule
\end{tabular}
\par\vspace{3pt}\parbox{0.95\linewidth}{\raggedright\scriptsize La neutralité coûte 1,6~Md€ cumulés sur 2027-2040 au
centre (0,6 à 2,1), soit 0,3~Md€ par an en régime de croisière. La règle (iii) coûte dès 2027, avant toute recette, et
donne aux héritiers qui ne vendent pas une baisse des droits ; sur le jeu bas, elle ne rapporte rien en cumul. La
règle (vi) rapporterait 6,4~Md€ par an dès 2028 au centre, mais taxe à la transmission. Le scénario prudent ne sert
qu'au bouclage : il reprend le jeu bas, avec deux ans de décalage, pour n'engager que des recettes constatées.}
\end{table}

\subsection{Pacte Dutreil, localisation des grands patrimoines, exit tax}

\emph{Pacte Dutreil} (art.~787~B : 75~\% de la valeur exonérée ; engagement individuel porté de 4 à 6~ans par la
LF 2026, sources secondaires [S]). La plus-value latente suit les titres et la correction porte sur les droits
réellement payés, qui sont faibles : $m$ vaut 5 à 10~\%. Entreprise de 10~M€ (prix 0,5~M€) transmise à deux enfants,
qui la cèdent après l'engagement : charge totale de 0,63~M€ aujourd'hui (6~\%), 3,34~M€ avec la réforme (33~\%), contre
5,58~M€ (56~\%) si le fondateur avait vendu avant de transmettre. Dans l'exemple à 1~M€ : 3~\% aujourd'hui, 30~\% avec
la réforme, 41~\% en vendant avant. L'écart restant n'est pas une interaction des deux impôts : c'est l'exonération
Dutreil elle-même, conservée par l'héritier qui vend après l'engagement. Si les héritiers gardent l'entreprise, rien ne
change. La hausse de charge pour les familles qui cèdent est forte, et elle est voulue [C] : le cumul « 75~\%
d'exonération + effacement » est supprimé, pas le pacte.

\emph{Localisation.} Une fois l'effacement fermé, la sortie qui reste est le départ. L'exit tax (art.~167~bis) impose
la plus-value latente au départ (participation d'au moins 50~\% ou titres de plus de 800\,000~€), avec sursis, puis
dégrève l'impôt si les titres sont encore détenus deux ans après le départ, cinq ans au-delà de 2,57~M€ [S : sources
secondaires, texte non relu]. La plus-value reprise entre dans sa base. Un héritier déjà non résident qui vend des
titres de portefeuille échappe en général à l'impôt français (conventions non vérifiées). La littérature dit trois
choses [E]. (1) Les très riches bougent : un point d'impôt sur la fortune en plus réduit de 2~\% environ le nombre de
contribuables fortunés en Scandinavie (Jakobsen, Kleven, Kolsrud, Landais et Muñoz, NBER WP 32153, \emph{AER} 2026) ;
après 2011, Madrid, seule région espagnole sans impôt sur la fortune, gagne 7,5~\% de contribuables fortunés en six ans
(Agrawal, Foremny et Martínez-Toledano, \emph{AEJ: Applied} 2025) ; 35~\% des milliardaires du Forbes 400 quittent
les États américains qui ont gardé un impôt successoral après 2001 (Moretti et Wilson, \emph{AEJ: Policy} 2023).
(2) Les recettes résistent : à Madrid comme aux États-Unis, l'impôt rapporte plus qu'il ne fait perdre, et en
Scandinavie 0,22~\$ seulement sur chaque dollar levé est perdu par départs. (3) En France, la réaction mesurée est
faible au niveau du dernier centile (0,02 à 0,23~\% des foyers de plus par point d'imposition effective, CAE,
Focus n°~118, 2025) ; les retours de redevables de l'IFI dépassent les départs depuis 2018 (France Stratégie, 2023).
Rien ne mesure le comportement du dernier dix-millième, qui porte l'essentiel de l'impôt latent. Ces élasticités
portent sur des impôts annuels sur le stock et ne se transposent pas telles quelles à un impôt payé une fois, à la
vente. Nous ne les appliquons pas [choix de méthode] : la prudence passe par le scénario de recettes prudent et par
deux mesures de garde [C] : délai de dégrèvement de l'exit tax porté à quinze ans pour les participations de plus de
2,57~M€ (durée en vigueur de 2014 à 2018, sources secondaires), et imposition, dans les limites des conventions, de
la plus-value transmise à un bénéficiaire non résident (à expertiser).

\subsection{Taux du PFU, CEHR et CDHR (inchangés depuis la version 2)}

Un point de PFU vaut 0,67~Md€ à comportements inchangés sur l'assiette des revenus de 2025 (67,0~Md€), 0,44~Md€ les
premières années avec l'élasticité de 0,77 de Lefebvre, Lehmann et Sicsic (2025) [S]. Une baisse à 25~\% coûterait
4,3~Md€ en mécanique, 3,1~Md€ après réaction. Les évaluations du PFU de 2018 ne trouvent d'effet que sur les dividendes
distribués (+15~\%, note IPP n°~46, 2019), aucun sur l'investissement (France Stratégie, 2023 ; Yagan, 2015, \emph{AER}) ;
Boissel et Matray (2022) ont été rétractés (\emph{AER} 113(7), 2023). CEHR et CDHR sont maintenues sur le capital
(coût nul) ; les supprimer coûterait 0,7 à 1,1 et 0,6~Md€. La CRDS s'éteint fin 2033 : le PFU passera à 30,9~\% sans
mesure nouvelle. L'indexation des plus-values sur l'inflation (0,6 à 0,9~Md€ par an) reste une orientation de long
terme.


\section{Le bouclage budgétaire, année par année}\label{ann:bouclage}
% Écrit par note8p/calculs/annexe_bouclage.py ; ne pas modifier à la main.
\subsection{Règles et calendrier}

Le bouclage tient trois règles. Aucune année de déficit organisée : le solde des quatre réformes reste compris entre
$-0{,}3$ et $+1{,}1$~Md€ chaque année. Les recettes de la fin de l'effacement des plus-values ne financent rien avant
d'être constatées ; elles forment une réserve. Les évolutions du droit en vigueur (fin de la contribution exceptionnelle
sur les bénéfices des grandes entreprises, $-7$~Md€ par an après l'exercice 2026 ; extinction de la CRDS fin 2033, qui
permet de baisser tous les taux du nouvel impôt de 0{,}54 point à recette constante) ne sont ni comptées ni utilisées.

Unités : euros constants de 2027. Les séries de l'IS, nominales (+3~\% par an), sont déflatées par leur propre
croissance ; les séries foncières et des plus-values sont au niveau de l'activité de 2025. L'écart de convention
(environ 4~\%) est sans effet sur les conclusions. L'IS d'un exercice est encaissé l'année suivante.

Calendrier [C] : nouvel impôt sur le revenu monté en charge sur quatre ans (2027-2030), chaque foyer payant
$(1-\lambda)\times$ ancien impôt $+\ \lambda\times$ nouvel impôt, $\lambda = 0{,}25$ à 1, ce qui garde la recette chaque
année ; taxe foncière nationale par tiers en 2029, 2030 et 2031, droits de mutation baissés le 1\textsuperscript{er}
janvier des mêmes années, avec clause de suspension ; taux réduit d'IS à 20~\% pour l'exercice 2027, aligné dès 2028 ;
déduction immédiate à 25~\% pour l'exercice 2028, 50~\% dès 2029 ; taux normal baissé par demi-points, au plus fort
possible sans qu'aucune année ultérieure passe sous $-0{,}3$~Md€ : 24~\% (exercice 2030), 23{,}5~\% (2032), 23~\%
(2033), 22{,}5~\% (2035). Étape conditionnelle : un point de plus dès l'exercice 2035 si la réserve atteint son niveau
central. Programme : \texttt{note8p/calculs/bouclage.py}.

\begin{table}[!htbp]\centering\scriptsize\setlength{\tabcolsep}{2.4pt}
\caption{Trajectoire complète (Md€ de 2027, années d'encaissement ; + = recette publique en plus)}\label{tab:traj_complete}
\begin{tabular}{@{}lrrrrrrrrrrrrrr@{}}\toprule
 & 27 & 28 & 29 & 30 & 31 & 32 & 33 & 34 & 35 & 36 & 37 & 38 & 39 & 40 \\ \midrule
Nouvel impôt sur le revenu (net) & 0{,}0 & 0{,}0 & 0{,}0 & $-$0{,}1 & $-$0{,}1 & $-$0{,}1 & $-$0{,}1 & $-$0{,}1 & $-$0{,}1 & $-$0{,}1 & $-$0{,}1 & $-$0{,}1 & $-$0{,}1 & $-$0{,}1 \\
Foncier net (taxe, droits, report, déductibilité) & 0{,}0 & $-$0{,}5 & 2{,}1 & 3{,}8 & 5{,}8 & 5{,}2 & 5{,}2 & 5{,}3 & 5{,}3 & 5{,}4 & 5{,}4 & 5{,}4 & 5{,}4 & 5{,}5 \\
Fin du taux réduit d'IS & 0{,}0 & 1{,}3 & 2{,}7 & 2{,}7 & 2{,}7 & 2{,}7 & 2{,}7 & 2{,}7 & 2{,}7 & 2{,}7 & 2{,}7 & 2{,}7 & 2{,}7 & 2{,}7 \\
Déduction immédiate de 50~\% & 0{,}0 & 0{,}0 & $-$3{,}7 & $-$6{,}5 & $-$5{,}1 & $-$3{,}9 & $-$3{,}0 & $-$2{,}1 & $-$1{,}3 & $-$0{,}8 & $-$0{,}7 & $-$0{,}7 & $-$0{,}7 & $-$0{,}7 \\
Baisse du taux normal & 0{,}0 & 0{,}0 & 0{,}0 & 0{,}0 & $-$2{,}8 & $-$2{,}8 & $-$4{,}3 & $-$5{,}8 & $-$5{,}9 & $-$7{,}4 & $-$7{,}4 & $-$7{,}4 & $-$7{,}4 & $-$7{,}4 \\
\textbf{Solde} & 0{,}0 & 0{,}8 & 1{,}0 & $-$0{,}1 & 0{,}6 & 1{,}1 & 0{,}6 & 0{,}0 & 0{,}7 & $-$0{,}2 & $-$0{,}1 & 0{,}0 & 0{,}0 & 0{,}0 \\
Réserve : fin de l'effacement (central) & 0{,}0 & 0{,}9 & 1{,}9 & 2{,}1 & 2{,}3 & 2{,}5 & 2{,}6 & 2{,}8 & 3{,}0 & 3{,}1 & 3{,}3 & 3{,}4 & 3{,}5 & 3{,}7 \\
Étape conditionnelle : baisse financée par la réserve & 0{,}0 & 0{,}0 & 0{,}0 & 0{,}0 & 0{,}0 & 0{,}0 & 0{,}0 & 0{,}0 & 0{,}0 & $-$3{,}0 & $-$3{,}0 & $-$3{,}0 & $-$3{,}0 & $-$3{,}0 \\
Solde avec la réserve et l'étape conditionnelle & 0{,}0 & 1{,}7 & 3{,}0 & 2{,}0 & 2{,}8 & 3{,}5 & 3{,}2 & 2{,}8 & 3{,}7 & $-$0{,}1 & 0{,}2 & 0{,}4 & 0{,}5 & 0{,}7 \\
Taux normal de l'IS (exercice précédent, \%) & 25{,}0 & 25{,}0 & 25{,}0 & 25{,}0 & 24{,}0 & 24{,}0 & 23{,}5 & 23{,}0 & 23{,}0 & 22{,}5 & 22{,}5 & 22{,}5 & 22{,}5 & 22{,}5 \\
Idem avec l'étape conditionnelle (\%) & 25{,}0 & 25{,}0 & 25{,}0 & 25{,}0 & 24{,}0 & 24{,}0 & 23{,}5 & 23{,}0 & 23{,}0 & 21{,}5 & 21{,}5 & 21{,}5 & 21{,}5 & 21{,}5 \\
\bottomrule\end{tabular}
\end{table}

\subsection{Sensibilités}

\textbf{Plus-values.} Au bas de la fourchette (1{,}0~Md€ en 2040 au lieu de 3{,}7), l'étape conditionnelle n'a pas
lieu ; rien d'autre ne change, puisque la réserve ne finance rien d'autre. \textbf{Foncier.} Avec une enveloppe de 10~Md€ au lieu de 15 (excédent net d'environ 1~Md€ par an), la baisse du
taux normal se limite à un point, et le calendrier de la déduction immédiate n'est plus finançable : le solde tombe à
$-3{,}2$~Md€ l'année du pic, sauf à retarder la déduction. Avec 20~Md€ (excédent de 10~Md€), le taux normal
descendrait à 21~\% dès 2035. Si l'évaluation des terrains prend un an de retard, la clause de suspension décale la
taxe et la baisse des droits, mais pas les mesures d'IS : sans autre précaution, le solde tomberait à $-1{,}8$~Md€ une
année. La déduction immédiate et chaque palier de baisse du taux doivent donc être soumis à la même clause que les
droits de mutation (programme : \texttt{FONCIER\_SCEN=E10|E20|E15\_glissement\_un\_an python3 calculs/bouclage.py}).
\textbf{IS.} La part imposable de l'investissement (75~\%, hypothèse) fait varier le coût de la déduction
dans la même proportion (pic de 5{,}4 à 8{,}1~Md€ pour 60 à 90~\%) ; inclure les véhicules utilitaires le multiplie
par 1{,}25. Le rendement de la fin du taux réduit va de 2{,}15 à 3{,}18~Md€. Dans le cas défavorable, la baisse du
taux normal doit attendre ; le calendrier par paliers le permet.


\section{Références principales}\label{ann:refs}
{\small\setlength{\parskip}{2pt}
Banzhaf, H. S. et N. Lavery (2010), « Can the land tax help curb urban sprawl? Evidence from growth patterns in
Pennsylvania », \emph{Journal of Urban Economics} 67(2). \par
Bérard, G. et A. Trannoy (2018), « The impact of the 2014 increase in the real estate transfer taxes on the French
housing market », \emph{Économie et Statistique} 500-502. \par
Bozio, A. et É. Wasmer (2024), \emph{Les politiques d'exonérations de cotisations sociales : une inflexion
nécessaire}, rapport au Premier ministre. \par
Chetty, R. (2009), « Is the taxable income elasticity sufficient to calculate deadweight loss? », \emph{AEJ: Economic
Policy} 1(2). \par
Chetty, R., A. Guren, D. Manoli et A. Weber (2013), « Does indivisible labor explain the difference between micro and
macro elasticities? », \emph{NBER Macroeconomics Annual} 27. \par
Chodorow-Reich, G., O. Zidar et E. Zwick (2024), « Lessons from the biggest business tax cut in US history »,
\emph{Journal of Economic Perspectives} 38(3). \par
Eerola, E., O. Harjunen, T. Lyytikäinen et T. Saarimaa (2021), « Revisiting the effects of housing transfer taxes »,
\emph{Journal of Urban Economics} 124. \par
France Stratégie (2023), Comité d'évaluation des réformes de la fiscalité du capital, rapport final. \par
Harberger, A. (1964), « The measurement of waste », \emph{American Economic Review} 54(3). \par
Hilber, C. et T. Lyytikäinen (2017), « Transfer taxes and household mobility », \emph{Journal of Urban Economics} 101.
\par
Hines, J. (1999), « Three sides of Harberger triangles », \emph{Journal of Economic Perspectives} 13(2). \par
Lehmann, É., F. Marical et L. Rioux (2013), « Labor income responds differently to income-tax and payroll-tax reforms »,
\emph{Journal of Public Economics} 99. \par
Maffini, G., J. Xing et M. Devereux (2019), « The impact of investment incentives: evidence from UK corporation tax
returns », \emph{AEJ: Economic Policy} 11(3). \par
Oates, W. et R. Schwab (1997), « The impact of urban land taxation: the Pittsburgh experience », \emph{National Tax
Journal} 50(1). \par
Zwick, E. et J. Mahon (2017), « Tax policy and heterogeneous investment behavior », \emph{American Economic Review}
107(1). \par
Les références propres à chaque bloc, avec ce qui en a été lu (texte ou résumé), figurent dans les comptes rendus
\texttt{note8p/sorties/\{ircsg,foncier,is,pfu\}.md}. Boissel et Matray (2022, \emph{AER}), souvent cité sur la hausse de
la fiscalité des dividendes de 2013, a été rétracté par ses auteurs en 2023 ; nous ne nous appuyons pas sur lui.\par}

\end{document}
